\section{An estimation assistant must preserve the target it validates}
\label{exec14:estimation-validation}

The choice of an estimator follows three earlier decisions: the business target,
the information available at the decision time, and identification under the
stated model. The assistant should keep these decisions separate from statistical
calibration and diagnostic adequacy. A balanced weighted sample can still omit
a common cause; an excellent predictive model can use a field recorded after
the intervention. A critic that approves the prose does not validate either
assumption. The causal co-pilot roadmap is valuable precisely because it makes
eligibility, baseline time, incident versus prevalent treatment, follow-up and
outcome definitions explicit before estimation \citep{v12_la_2407_19118}.

For an early arrears-review offer, let $X$ be the information available before
an eligible account's first qualifying arrears date, $A\in\{0,1\}$ the offer,
and $Y$ discounted net recovery over 180 days. The offer, acceptance and eventual
restructuring terms are distinct actions. The population $P$ consists of the
specified eligible accounts at time zero, with one declared interference and
sampling unit. Define $\tau_P=E_P\{Y(1)-Y(0)\}$. Offer cost is either included
in $Y$ or subtracted once when comparing actions. Changing the outcome horizon
or the eligible population creates a different target.

Under consistency, conditional exchangeability, positivity, a suitable sampling
law and the declared measurement conditions, write
$\mu_a(x)=E(Y\mid A=a,X=x)$ and $p(x)=\Pr(A=1\mid X=x)$. For candidate nuisance
functions $m_a$ and $e$, a population augmented inverse-probability score is
\[
 \phi=m_1(X)-m_0(X)+\frac{A\{Y-m_1(X)\}}{e(X)}
       -\frac{(1-A)\{Y-m_0(X)\}}{1-e(X)}.
\]
Conditioning on $X$ and subtracting the identified conditional contrast gives
\begin{equation}
 E(\phi\mid X)-(\mu_1-\mu_0)
 =(e-p)\left\{\frac{m_1-\mu_1}{e}
               +\frac{m_0-\mu_0}{1-e}\right\}.
 \label{exec14:aipw-remainder}
\end{equation}
This algebra explains the population double-robustness property when the
denominators are well defined: the conditional bias vanishes if the propensity
is right or both outcome regressions are right. It does not establish the causal
assumptions, a finite-sample confidence guarantee, or arbitrary learner validity.
Cross-fitting can separate nuisance training from score evaluation, but rates,
moments, dependence and uncertainty calculations still require their own
conditions. The inspected OCI notebook uses a cross-fitted AIPW approach within
a human-principal and actor/critic workflow; its source results remain distinct
from this project derivation \citep{v12_la_2607_22443,exec14_estimation_sources}.

\subsection{Full, trimmed and overlap targets answer different questions}

The exact three-stratum fixture in Table~\ref{exec14:target-table} gives a
concrete test of target preservation. Its probabilities, effects and monetary
interpretation are stipulated; no bank or package was fitted. The strata have
masses $p_s$, offer probabilities $e_s$ and conditional effects $\tau_s$.
\begin{table}[htbp]
\centering
\caption{Declared population law for a target-preservation fixture.}
\label{exec14:target-table}
\begin{tabular}{lrrr}
\toprule
Stratum & $p_s$ & $e_s$ & $\tau_s$\\
\midrule
Low offer propensity & $1/2$ & $1/20$ & $0$\\
Intermediate propensity & $3/10$ & $1/2$ & $4$\\
High offer propensity & $1/5$ & $19/20$ & $12$\\
\bottomrule
\end{tabular}
\end{table}
The full-population contrast is $\sum_s p_s\tau_s=18/5=3.6$. A fixed rule that
retains only propensities in $[.1,.9]$ keeps the intermediate stratum and yields
$4$. Conventional overlap weighting instead gives
\[
 \tau_{\mathrm{ov}}=
 \frac{\sum_s p_s e_s(1-e_s)\tau_s}{\sum_s p_s e_s(1-e_s)}
 =\frac{1656}{433}\simeq3.82448.
\]
These are three correct answers to three different questions. The target rule,
retained mass and weighting measure belong in the report next to the number.
Stabilizing an estimate by removing low-overlap accounts may be scientifically
reasonable, but it changes which accounts the recommendation concerns.

The pinned OCI notebook supplies an instructive source-level discrepancy. It
trims a propensity range, then applies \texttt{aipw\_summary} to the retained
pseudo-outcomes with uniform averaging while serializing that result as ATO.
The operation is a trimmed-population average; it is not automatically the
conventional $e(X)\{1-e(X)\}$-weighted overlap target. The proposed adapter must
rename the target with its rule or implement the intended weighting explicitly.
When trimming uses estimated propensities, the retained population is itself
estimated, adding selection and target uncertainty. The exact fixture above
uses known propensities and establishes only the arithmetic distinction
\citep{exec14_estimation_sources,exec14_estimation_oracles}.

The paper reports 231 ACIC datasets, a ten-dataset scaffolding comparison and
disguised LaLonde analyses. It reports better behavior with its structured
workflow and critic in those settings; these findings justify a candidate
architecture, not a claim of independently reproduced bank performance. A
critic's overlap refusal may be useful even when no effect is returned.
However, selecting critic-approved estimates does not automatically establish
coverage conditional on approval. Both the all-request and selected denominators
must be shown, with the selection rule fixed before final evaluation
\citep{v12_la_2607_22443}.

\subsection{Validation losses, unknown states and retrieval failure}

CausalDS separates narrative scenes, an observation layer and hidden causal
truth. Its identification targets refer to conceptual variables; an adapter must
not infer that a noisy observed proxy provides the same identifying law. The
paper's 100-scene exam, sampled from 953 scenes, reports different patterns for
symbolic success, numeric quality, abstention and resource cost. Those are dated
author findings for a benchmark. Its source scorer also exposes implementation
choices that matter when adapting the benchmark: missing identification metadata
can default toward an identifiable interpretation, and finite numeric endpoints
must be checked separately from successful conversion to a floating-point value
\citep{v12_la_2607_08093,exec14_estimation_sources}.

The interval score for a central interval $[l,u]$, target $\theta$ and nominal
error $\alpha$ is
\[
 S_\alpha(l,u;\theta)=(u-l)+\frac{2}{\alpha}(l-\theta)\mathbf1\{\theta<l\}
                  +\frac{2}{\alpha}(\theta-u)\mathbf1\{\theta>u\}.
\]
It penalizes width and noncoverage. The inspected CausalDS source reports both
uncapped and capped normalized scores, together with coverage and width. It is
incorrect to describe it as offering only capped loss. A dashboard that keeps
only the capped value can nevertheless reverse a comparison. Score pairs
$(1,100)$ and $(9,9)$ have uncapped means $50.5$ and $9$; with a cap of ten,
the means are $5.5$ and $9$. Such pairs can be constructed with covered
zero-centered intervals of the corresponding widths and truth zero, so the
arithmetic is compatible with the interval-score definition. This is an exact
illustration of a dashboard choice, not a result from executing CausalDS
\citep{exec14_estimation_sources,exec14_estimation_oracles}.

The proposed assessment keeps identified, nonidentified and unknown as separate
scientific states. Unsupported operations, malformed files, execution failures
and deliberate abstentions have separate operational states. It reports how
many requests returned finite valid intervals and how many were refused, while
retaining uncapped losses for the valid interval population. Assigning an
infinite or arbitrary loss to every refusal may be inappropriate for a task
where refusal is correct; silently dropping failures is also inappropriate.
Task-specific scientific scoring and the all-request operating denominator
must therefore be published together.

ARGUS provides a complementary example at the assumption-audit boundary. It
audits reported support for identification assumptions, not causal truth. Its
synthetic flaw results show that retrieval gating changes the tradeoff between
flagging flaws and missing omissions. In its five-paper human pilot, weak
retrieval produces excessive severity, and the reported adjustment is in-sample.
The useful remedy is to retain retrieval quality and evidence status as separate
fields, with unresolved evidence labeled unknown. A failure to retrieve a bank
policy is not evidence that the policy does not exist
\citep{v12_la_2609_30867}.

The initial prototype should adapt OCI's specification and critic artifacts
behind fixed bounded operations, and adapt CausalDS's hidden-truth and scoring
separation only after its generator and scoring paths are independently checked.
OCI's unbounded notebook-runner timeout is not an admissible default for the
bounded CLI/MCP service. Direct package integration, a live provider comparison
and bank validation are separate later claims. The concrete near-term output is
a reproducible design that preserves the target, refuses unsupported certainty,
and measures review burden against a deterministic analyst baseline.
\subsection{A workflow succeeds by preserving a decision}

The developed examples give each handoff a substantive test. A servicing rule
must react to its generated history; a rollout comparison must preserve cohort
composition and untreated-trend assumptions; an observation correction must
retain the population frame and measurement law. The architectural examples
then show how trimming, repeated criticism, failed tools and selective scoring
can change what a completed report means. These failures require different
remedies and cannot be summarized by one success rate.

The proposed controller should preserve the question, unresolved premises and
evidence through those operations, with deterministic checks wherever their
conditions can be made explicit. Its advantage over a structured analyst
workflow remains an empirical question about usable answers and complete
review cost. The next chapter compares the tools that could implement this
restricted role, including the source, rights and execution obligations that
a plausible architecture alone cannot settle.
\label{exec14:estimation-validation-end}
